% 灰度打印版颜色覆盖 — 由 Makefile 在 \printversiontrue 时 \input
% 填充色 → 白/浅灰（保证文字可读）
\colorlet{bglight}{white}
\colorlet{lightblue}{gray!15}
\colorlet{skyblue}{gray!15}
\colorlet{inputyellow}{gray!10}
\colorlet{outputpink}{gray!20}
\colorlet{lightpink}{gray!20}
\colorlet{lightgreen}{gray!15}
\colorlet{lightorange}{gray!25}
% 强调/边框色 → 黑/深灰
\colorlet{deepblue}{black}
\colorlet{deeppink}{gray!35}
\colorlet{dangerred}{gray!55}
\colorlet{vibrantred}{gray!55}
\colorlet{improvedpink}{gray!40}
% 曲线色 — deepblue(Baseline)=黑, vibrantgreen(MIKU)=深灰, 确保可区分
\colorlet{vibrantgreen}{gray!45}
\colorlet{vibrantorange}{gray!65}
% 别名跟随
\colorlet{obsstatic}{gray!50}
\colorlet{bandok}{gray!15}
\colorlet{obsbuf}{gray!50}
\colorlet{egopt}{black}
